\documentclass[12pt]{article}
\usepackage[T1]{fontenc}
\usepackage[utf8]{inputenc}
\usepackage[brazil]{babel}
\usepackage[margin=1in]{geometry}
\setlength{\parindent}{0pt}
\begin{document}
\section*{Tema 16}
Processo: PEDILEF 0501999-48.2009.4.05.8500/ SE\\
Situacao: Julgado\\
Ramo: DIREITO ADMINISTRATIVO\\
Questao: Saber qual o termo inicial da progressão funcional de policial federal.\\
Tese: A eficácia da progressão funcional deve ser observada segundo a situação individual de cada servidor e seus efeitos retroagem ao momento em que os requisitos legais foram implementados. Vide Tema 82 e PEDILEF n. 500367784.2014.404.7101. Entendimento revisado pelo PEDILEF n. 05207128420124058300 (orientação alinhada ao STJ: RESP 1.649.269/RJ).\\
Fonte: https://www.cjf.jus.br/phpdoc/virtus/
\end{document}
